% 02_introduction.tex
\section{Introduction}
Learning Data Structures and Algorithms (DSA) is a fundamental stepping stone for computer science students and software engineers. However, the abstract nature of these concepts---such as tree balancing rotations, dynamic graph traversals, and heap properties---often poses a significant challenge. Traditional learning methods relying on static textbook diagrams or raw console outputs frequently fail to capture the dynamic, step-by-step transformations that occur in memory during algorithm execution.

To bridge this educational gap, the \textbf{Data Structure Visualizer} was developed. It is a high-performance, interactive, and visually engaging desktop application designed to bring abstract algorithms to life. Built using modern C++, the Simple and Fast Multimedia Library (SFML), and Dear ImGui, this application transcends standard educational tools by providing a seamless, game-like user experience.

\subsection{Key Features}
The application stands out by offering a comprehensive suite of interactive features designed to enhance cognitive understanding:
\begin{itemize}[leftmargin=*]
    \item \textbf{Timeline Playback System:} Users can play, pause, step forward, or step backward through an algorithm's execution. Operating like a ``movie projector'' for data structures, it allows users to deeply analyze complex steps at their own pace, with fully adjustable execution speeds.
    \item \textbf{Real-time Pseudocode Highlighting:} As the visualizer animates a specific step, the corresponding line of pseudocode is highlighted in a dedicated Code Panel. This creates a direct cognitive link between abstract code and visual behavior.
    \item \textbf{Cinematic UI/UX and Animations:} Moving away from rigid, snapping graphics, the application utilizes linear interpolation (Lerp), spring physics, and color blending. Every node movement, tree rotation, and state transition is rendered with extreme smoothness to prevent visual confusion.
    \item \textbf{Dynamic Theming:} Features a meticulously designed interface with full support for Light Mode and Dark Mode. The user interface integrates Cyberpunk and Holographic aesthetics to keep the learning process visually stimulating.
    \item \textbf{Flexible Data Input:} Users can initialize and modify structures via manual keyboard input, randomized generation, or by importing external datasets from files.
\end{itemize}

\subsection{Supported Data Structures and Algorithms}
The visualizer currently supports a robust library of standard and advanced data structures:
\begin{itemize}[leftmargin=*]
    \item \textbf{Linear Structures:} Singly Linked List (Insert, Delete, Search, Update).
    \item \textbf{Trees \& Heaps:} AVL Tree (featuring automatic self-balancing rotations), Min Heap, and Max Heap.
    \item \textbf{Grid Environments:} Support for maze generation, obstacle placement, Breadth-First Search (BFS), and the A* Pathfinding algorithm.
    \item \textbf{Node-Edge Graphs:} Directed and Undirected graphs represented via Adjacency Matrices or Adjacency Lists. Supported algorithms include Dijkstra, Bellman-Ford, Floyd-Warshall, Johnson's Algorithm, and A*.
\end{itemize}

\subsection{Technical Architecture}
Under the hood, the application follows a strict \textit{Model-View-Controller (MVC)} architecture. 
\begin{itemize}
    \item \textbf{The Model} handles the mathematical data representations and generates a timeline of discrete \textit{Frames}.
    \item \textbf{The View} is powered by SFML for high-framerate 2D hardware-accelerated rendering and Dear ImGui for creating flat-design, highly responsive control panels. 
    \item \textbf{The Controller} acts as the central engine (\textit{AppEngine}), bridging user inputs, playback states, and continuous visual updates.
\end{itemize}

By combining robust algorithmic logic with modern visual aesthetics, the Data Structure Visualizer serves as both an effective self-study companion and a powerful demonstration tool for computer science education.
